% ===== tr-06 块1：第 4 章 pp.52--56（原文页 52--56）=====
% [边界] 本文件首句接 tr-05（tr05f）p51 末句 “These rays tell us that …”，属同一自然段，衔接处不要插空行。
% [原文 p.52]
此时，我们会得到不同的自同构，取决于我们是在奇点 $p$ 的上方还是下方穿过。若从下方穿过，射线告诉我们取 $R^k_{\tau,\tau,\sigma_1}$ 的一个自同构，它由
\begin{equation}
x\mapsto x(1+w),\quad y\mapsto y(1+w)^{-1},\quad w\mapsto w, \tag{4.10.2}
\end{equation}
给出；而若从奇点上方穿过，则使用自同构
\begin{equation}
x\mapsto x(1+w^{-1}),\quad y\mapsto y(1+w^{-1})^{-1},\quad w\mapsto w. \tag{4.10.3}
\end{equation}

其实要注意，$1+w$ 或 $1+w^{-1}$ 在 $R^k_{\tau,\tau,\sigma_i}$ 中并不可逆，因此我们需要把该环在 $1+w$ 处（等价地，在 $1+w^{-1}$ 处）局部化，以作修改。下面来看这会如何影响纤维积 $R^k_{\tau,\sigma_1,\sigma_1}\times_{R^k_{\tau,\tau,\sigma_2}}R^k_{\tau,\sigma_2,\sigma_2}$。
% [Tutor Note: 式 (4.10.2) 与 (4.10.3) 并不独立：$w$ 在该环中可逆，且 $1+w^{-1}=w^{-1}(1+w)$，故复合 (10.3)$\circ$(10.1) 得 $x\mapsto xw(1+w^{-1})=x(1+w)$、$y\mapsto yw^{-1}(1+w^{-1})^{-1}=y(1+w)^{-1}$，恰为 (10.2)；这正是下文“两种走法给出同一纤维积”的计算依据。]

若采用奇点下方的平行移动，则映射 $R^k_{\tau,\sigma_1,\sigma_1}\to R^k_{\tau,\tau,\sigma_2}$ 恰好由 (10.2) 给出，而 $R^k_{\tau,\sigma_2,\sigma_2}\to R^k_{\tau,\tau,\sigma_2}$ 仍取典范的那个。此时可得
\[
R^k_{\tau,\sigma_1,\sigma_1}\times_{R^k_{\tau,\tau,\sigma_2}}R^k_{\tau,\sigma_2,\sigma_2}\cong\C[X,Y,W^{\pm},t]/(XY-(1+W)t,t^{k+1}),
\]
其中
\[
X=(x,x(1+w)),\quad Y=(y(1+w),y),\quad W=(w,w),\quad t=(xy,xy).
\]
% [Tutor Note: 上式环中的 $W^{\pm}$ 是 $W^{\pm 1}$ 的简写（同页上文与 p.51 的同类式子均写作 $W^{\pm 1}$），不是 $W$ 的 $\pm$ 次幂之外的记号。]

另一方面，若采用奇点上方的平行移动，我们需要把自同构 (10.3) 与同构 (10.1) 复合，得到一个由下式给出的映射 $R^k_{\tau,\sigma_1,\sigma_1}\to R^k_{\tau,\tau,\sigma_2}$：
\[
x\mapsto xw(1+w^{-1})=x(1+w),\quad y\mapsto yw^{-1}(1+w^{-1})^{-1}=y(1+w)^{-1},\quad w\mapsto w.
\]

于是该映射与 (10.2) 完全相同，从而我们得到同一个纤维积。粘合所得的加厚与选择无关。额外引入的这些自同构消去了单值化所带来的问题。

这是一个非常局部的现象。随之而来的下一个问题是：在更全局的情形下，我们需要把附着于射线的自同构传播出去。事实上，可以设想我们所看的图景处在一个更复杂的局面之中，正如 图 4.7 左侧所示。那里有两个奇点，且从每个奇点向各个方向都有射线发出。让我们遵循如下规则：任何涉及穿过某条射线的平行移动的环的同构，都必须用上面描述的自同构作相应修正。于是，比如说看顶点 $v_1$，我们需要把五个不可约分支粘合起来，但其中只有一个粘合被修正。这样得到的粘合不会相容。为修正这一点，可以把射线无限延伸，并沿射线“平行移动”该自同构。这件事有一种确切的做法。图 4.7 右图显示了这一做法，虚线表示射线的延伸。现在，若沿一个方向穿过一条射线所产生的是沿另一方向穿过该射线所得自同构的逆，
% [原文 p.53]
\figimg[0.72]{p4_fig07.png}
\figtitle{4.7}{}
% [Tutor Note: 原文图注只有“Figure 7.”，无描述性文字；图本身含左右两幅图：左图两端点 $v_1$、$v_2$ 各有射线发出，右图以虚线画出射线的延伸。译文图注留空，以忠实原文。]

\figimg[0.72]{p4_fig08.png}
\figtitle{4.8}{}
% [Tutor Note: 原文图注只有“Figure 8.”，无描述性文字；图为从两射线交点向各方向新增（虚线）射线后的散射图，对应正文“添加从旧射线交点发出的新射线”。]

则可以看到，顶点 $v_1$ 与 $v_2$ 处的粘合现在变得相容了。

然而，在两射线的交点处又会冒出一个新问题。同样地，当我们用平行移动和由穿过射线所诱导的自同构来认同各个环时，我们不希望这种选择依赖于所走的具体路径。由于一般而言附着于两条射线的两个自同构并不交换，我们在交点处又会遇到麻烦。

事实上，这正是 2004 年初、Kontsevich 与 Soibelman 的论文 [54] 出现前不久，我们的思考所处的位置。这一问题的解决办法------实际上也是 Kontsevich 与 Soibelman 论证的关键部分------是添加从旧射线的交点发出的新射线，如 图 4.8 所示。这些射线的添加方式保证：绕交点走一圈所给出的自同构的复合实际上是恒等，从而这些认同将与路径的选取无关。

这里的描述有些含糊，但展示了基本想法。我们已经看到如何通过把一些基本片段粘合起来而得到退化。除了这些不同的基本片段之外，在二维情形下，我们的方法与 Kontsevich 与 Soibelman 在 [54] 中所采用的方法的主要区别在于：我们是在与 [54] 所采用的那个仿射结构相对偶的仿射结构中工作。
% [原文 p.54]
他们沿梯度流线传播自同构，而我们能够沿关于仿射结构的直线传播自同构。这在更高维情形省去了大量麻烦，因为在那些地方梯度流线要难控制得多。这使我们能够在所有维数上得到结果。

当然，我们尚未把究竟如何编码自同构、以及它们如何传播讲得足够清楚，但至少第一点我们会在下一节讲得更明白。至于第二点，关键在于它们沿直线传播；这一点实际上对于保证自同构不会开始涉及在 $X_0$ 的不可约分支上带极点的单项式至关重要。所以这里我们已经看到了某种看起来已经是热带的东西：射线的并看起来像一棵热带树。同样地，我们会在下一节把这一点讲得更精确。

在更高维情形，论证要微妙得多。不再是射线携带自同构，而是余维一的墙携带自同构，并且必须非常小心这些墙如何传播。此外，在判别轨迹附近，算法的收敛性还有很大的技术困难。在 [54] 中，二维情形是通过如下论证处理的：可以保证新加入的射线避开每个奇点的某个邻域，但这依赖于对度量的仔细选取。在更高维情形这不再成立，我们改用代数方法来证明收敛性。所有这些困难都在 [33] 中得到了克服。

在 [25] 中我写出了二维情形证明的完整版本；其优点是避开了大部分真正技术性的问题。希望 [25] 能比主论文 [33] 为这里概述的想法提供一个更平易的入口。

现在我们转向对所涉及的自同构作更精确的描述，并给出证据说明：对显式变形的描述（我们视之为 B-模型信息）确实编码了镜像上的 A-模型信息。

\section{热带顶点}

为使讨论简化，本节我们只用上一节中出现过的最简单的环，即形如 $R^k_{\sigma,\sigma,\sigma}$ 的环，其中 $\sigma$ 是一个极大（二维）胞腔。该环同构于 $\C[x^{\pm1},y^{\pm1},t]/(t^{k+1})$。我们改为在形式框架下工作，令
\[
R=\C[x^{\pm1},y^{\pm1}][\![t]\!].
\]
这是以 $x$、$y$ 的 Laurent 多项式为系数的 $t$ 的形式幂级数环。设 $f\in R$ 具有形式
\[
f=1+tx^ay^b\cdot g(x^ay^b,t),\quad g(z,t)\in\C[z][\![t]\!].
\]
% [Tutor Note: 原文的形式幂级数环把 $t$ 的方括号排成黑板体（bbold 字体）的粗方括号，以区别于多项式环；本译本统一用 $[\![t]\!]$ 表示形式幂级数。注意 $R$ 是把 $t$ 形式化之后的环，故下文 $R^k_{\sigma,\sigma,\sigma}$ 与 $R$ 的关系是模 $t^{k+1}$ 截断。]
% [原文 p.55]
于是这定义了一个 $R$ 的 $\C[\![t]\!]$-代数自同构 $\theta_{(a,b),f}$，它由
\[
\theta_{(a,b),f}(x)=x\cdot f^b,\quad \theta_{(a,b),f}(y)=y\cdot f^{-a}
\]
给出。注意 $\theta_{(a,b),f}^{-1}=\theta_{(a,b),f^{-1}}$。这些自同构还具有如下性质：它们保持全纯辛形式 $\frac{dx}{x}\wedge\frac{dy}{y}$。
% [Tutor Note: $\theta_{(a,b),f}^{-1}=\theta_{(a,b),f^{-1}}$ 可直接用定义验证：$\theta_{(a,b),f}$ 把 $x$ 乘 $f^b$、把 $y$ 乘 $f^{-a}$，而以 $f^{-1}$ 代 $f$ 即得逆映射。]

我们把\emph{热带顶点群} $\mathbb{V}$ 定义为：由所有这样的自同构生成的 $R$ 的 $\C[\![t]\!]$-代数自同构子群，关于极大理想 $(t)\subseteq\C[\![t]\!]$ 作完备化。注意在 $\mathbb{V}$ 中，无穷乘积只有在如下情形才有定义：对每个 $k>0$，只有有限多个因子模 $t^k$ 非平凡。这是对 Kontsevich 与 Soibelman 在 [54] 中最初引入的群的一个轻微改动。

现在我们描述上一节所述射线的一个局部版本。为方便起见，令 $M=\Z^2$、$M_\R=M\otimes_\Z\R$，并把 $\C[x^{\pm1},y^{\pm1}]$ 与 $\C[M]$ 等同。

\begin{definition}{4.11.1}
$M_\R$ 中的\emph{射线}或\emph{直线}是一个对 $(\mathfrak{d},f_\mathfrak{d})$，其中当 $\mathfrak{d}$ 是射线时 $\mathfrak{d}=\R_{\le0}m$，当 $\mathfrak{d}$ 是直线时 $\mathfrak{d}=\R m$，这里 $m\in M\setminus\{0\}$。此外，
\[
f_\mathfrak{d}=1+tz^m\cdot g(z^m,t)\in R,\quad g(z,t)\in\C[z][\![t]\!].
\]

一个\emph{散射图} $\mathfrak{D}$ 是一族射线与直线 $\{(\mathfrak{d},f_\mathfrak{d})\}$，具有性质：对任意 $k>0$，除有限多个元素外，其余元素都满足 $f_\mathfrak{d}\equiv1\bmod t^k$。

给定散射图 $\mathfrak{D}$，令
\[
\operatorname{Supp}\mathfrak{D}=\bigcup_{(\mathfrak{d},f_\mathfrak{d})\in\mathfrak{D}}\mathfrak{d}.
\]

若给定一条道路 $\gamma:[0,1]\to M_\R\setminus\{0\}$，满足 $\gamma(0),\gamma(1)\notin\operatorname{Supp}(\mathfrak{D})$，且 $\gamma$ 与它所穿过的每条射线都横截，则我们可以定义\emph{路径有序积} $\theta_{\gamma,\mathfrak{D}}\in\mathbb{V}$，它是与 $\gamma$ 所穿过的每条射线相对应的自同构的复合。我们对每个幂次 $k>0$ 分别定义路径有序积，然后取极限。对任意给定的 $k>0$，可以找到数
\[
0<t_1\le t_2\le\cdots\le t_s<1
\]
以及元素 $\mathfrak{d}_i\in\mathfrak{D}$，使得 $f_{\mathfrak{d}_i}\neq1\bmod t^k$，$\gamma(t_i)\in\mathfrak{d}_i$，且当 $t_i=t_j$ 时 $\mathfrak{d}_i\neq\mathfrak{d}_j$，并且 $s$ 取尽可能大。对每个 $i$，按下述方式定义自同构 $\theta_{\mathfrak{d}_i}$。设 $n\in N=\operatorname{Hom}(M,\Z)$ 是唯一的本原元素，它在 $\mathfrak{d}_i$ 上取零、在 $\gamma'(t_i)$ 上取负值。然后定义 $\theta_{\mathfrak{d}_i}$ 为自同构
\[
\theta_{\mathfrak{d}_i}(z^m)=z^m f_{\mathfrak{d}_i}^{\langle n,m\rangle}.
\]
注意这具有 $\theta_{(a,b),f_\mathfrak{d}}$ 的形式（对适当的 $(a,b)$ 选取）。于是我们定义
\[
\theta^k_{\gamma,\mathfrak{D}}=\theta_{\mathfrak{d}_s}\circ\cdots\circ\theta_{\mathfrak{d}_1}.
\]
\end{definition}
% [Tutor Note: 射线/直线的记号 $\mathfrak{d}$ 是欧拉 Fraktur 体的小写 d（页图字体 EUFM10），不是偏导符号 $\partial$；散射图用 Fraktur 体的 $\mathfrak{D}$。文字层把它们抽成普通的 d、D，无法据以区分，必须以页图为准。]
% [Tutor Note: “唯一一个本原元素”即 primitive 向量：$n\in N=\operatorname{Hom}(M,\Z)$ 在 $\mathfrak{d}_i$ 上为零（即与射线方向正交），在 $\gamma'(t_i)$ 上为负，用以保证穿过方向与自同构取值方向一致；$\langle n,m\rangle$ 是自然配对 $N\times M\to\Z$。]
% [原文 p.56]
注意，若 $\gamma$ 同时穿过两条射线，则顺序无关紧要，因为容易验证：若两个自同构对应同一个底空间集 $\mathfrak{d}\subseteq M_\R$，则它们交换。最后，我们定义
\[
\theta_{\gamma,\mathfrak{D}}=\lim_{k\to\infty}\theta^k_{\gamma,\mathfrak{D}}.
\]

于是我们可以在这一框架中表述 [54] 的本质引理：

\begin{proposition}{4.11.2}
设 $\mathfrak{D}$ 是一个散射图。则存在一个散射图 $S(\mathfrak{D})$，使得 $S(\mathfrak{D})\setminus\mathfrak{D}$ 只由射线组成，且对绕原点的任意闭路 $\gamma$，$\theta_{\gamma,S(\mathfrak{D})}$ 都是恒等映射。
\end{proposition}

这个证明非常简单，而且是算法性的；我给出一个简短的概述。人们构造一列散射图 $\mathfrak{D}_1=\mathfrak{D},\mathfrak{D}_2,\dots$，具有性质 $\theta_{\gamma,\mathfrak{D}_k}\equiv\mathrm{id}\bmod t^k$。对 $\mathfrak{D}_1$ 这显然成立，于是我们归纳地进行：假设已构造出 $\mathfrak{D}_k$。然后证明（通过考察 $\mathbb{V}$ 的李代数）：
\begin{gather}
\theta_{\gamma,\mathfrak{D}_k}(x)=x\sum_{i=1}^{n}b_ic_it^kx^{a_i}y^{b_i}\\
\theta_{\gamma,\mathfrak{D}_k}(y)=-y\sum_{i=1}^{n}a_ic_it^kx^{a_i}y^{b_i}
\end{gather}
其中 $a_i,b_i$ 为整数（$a_i,b_i$ 不全为零），$c_i\in\C$。
% [Tutor Note: 这里两式的等号应理解为模 $t^{k+1}$：由归纳假设 $\theta_{\gamma,\mathfrak{D}_k}\equiv\mathrm{id}\bmod t^k$，其与恒等的偏差恰为该 $t^k$ 阶项；作者在下一段给出的自同构处才明确写出 $\bmod t^{k+1}$。]

接着，通过添加射线
\[
(\R_{\le0}(a_i,b_i),\,1\pm c_it^kx^{a_i}y^{b_i}),\quad 1\le i\le n
\]
而得到 $\mathfrak{D}_{k+1}$，其中符号的选取使得当 $\gamma$ 穿过这条射线时，产生自同构
\[
x\mapsto x(1-b_ic_it^kx^{a_i}y^{b_i})\bmod t^{k+1},\quad y\mapsto y(1+a_ic_it^kx^{a_i}y^{b_i})\bmod t^{k+1}.
\]

由于该自同构将与 $\mathfrak{D}_k$ 中所有其他自同构在模 $t^{k+1}$ 的意义下交换，插入这些射线将恰好抵消掉对 $\theta_{\gamma,\mathfrak{D}_k}$ 的 $k$ 阶贡献，从而 $\theta_{\gamma,\mathfrak{D}_{k+1}}\equiv\mathrm{id}\bmod t^{k+1}$。

很容易把这个算法编程并探索这些散射图。它们似乎具有非常丰富而迷人的结构。下面的简单例子显示了它们的复杂性。

\begin{example}{4.11.3}
考虑如下情形：
\[
\mathfrak{D}=\{(\R(1,0),(1+tx^{-1})^\ell),\,(\R(0,1),(1+ty^{-1})^\ell)\}
\]
其中 $\ell$ 为某个正整数。对 $\ell=1$，容易验证
\[
S(\mathfrak{D})\setminus\mathfrak{D}=\{(\R_{\ge0}(1,1),\,1+t^2x^{-1}y^{-1})\}.
\]
图 4.9 明确展示了当沿所描绘的环道行进时这些自同构是什么；读者容易验证这五个自同构的复合是恒等映射。
\end{example}
% [Tutor Note: 例中的 $\ell$ 是斜体字母 ell，不是数字 1；$(1+tx^{-1})^\ell$ 的指数若从文字层转抄极易误作指数 1，须以页图为准。]
% [Tutor Note: $\R(1,0)$、$\R(0,1)$ 是直线（定义 4.11.1 中 $\mathfrak{d}=\R m$），不是射线；新增支撑 $\R_{\ge0}(1,1)$ 才是从原点出发指向第一象限的射线。又 $f=(1+tx^{-1})^\ell$ 中的 $x^{-1}=z^{m}$ 对应 $m=(-1,0)$，仍落在直线 $\R(1,0)$ 上。]
